\clearpage
\subsection{Evaluation metrics}
\label{sec:metrics}

Let $c_b$, $b=1,\ldots,B$, denote the test conditions.  A future trajectory
of length $L$ with $d$ state variables is represented by a vector in
$\mathbb{R}^{D}$, where $D=Ld$.  For condition $c_b$, let
\[
P_b := P(\,\cdot\mid c_b),
\qquad
Q_{\theta,b} := Q_\theta(\,\cdot\mid c_b)
\]
denote the reference and learned conditional distributions, respectively.
The model produces an ensemble
$\{\widehat{\mathbf y}_b^{(m)}\}_{m=1}^{M}$ from $Q_{\theta,b}$.
All reported metrics are averaged over test conditions, and lower values are
better.

\subsubsection{Deterministic ground truth}

For deterministic systems, $P_b=\delta_{\mathbf y_b}$ is concentrated at the
single ground-truth future $\mathbf y_b\in\mathbb{R}^{D}$.

\paragraph{Continuous ranked probability score (CRPS).}
CRPS is a proper scoring rule for comparing a scalar predictive cumulative
distribution function $F$ with an observation $y$
\citep{gneiting2007strictly}.  It is defined as
\begin{equation}
  \operatorname{CRPS}(F,y)
  =
  \int_{-\infty}^{\infty}
  \left[F(a)-\mathbf{1}\{a\geq y\}\right]^2\,\mathrm{d}a.
  \label{eq:crps-cdf}
\end{equation}
Equivalently, if $X$ and $X'$ are independent draws from $F$, then
\begin{equation}
  \operatorname{CRPS}(F,y)
  =
  \mathbb{E}|X-y|
  -\frac{1}{2}\mathbb{E}|X-X'|.
  \label{eq:crps-expectation}
\end{equation}
We apply this score to each scalar entry of the future trajectory.  Its
finite-ensemble estimate, averaged over conditions and trajectory entries, is
\begin{equation}
  \widehat{\operatorname{CRPS}}
  =
  \frac{1}{BD}\sum_{b=1}^{B}\sum_{j=1}^{D}
  \left[
    \frac{1}{M}\sum_{m=1}^{M}
    \left|\widehat y_{b,j}^{(m)}-y_{b,j}\right|
    -
    \frac{1}{2M^2}\sum_{m=1}^{M}\sum_{m'=1}^{M}
    \left|\widehat y_{b,j}^{(m)}-\widehat y_{b,j}^{(m')}\right|
  \right].
  \label{eq:crps-ensemble}
\end{equation}

\paragraph{Root mean squared error (RMSE) and mean absolute error (MAE).}
Both metrics use the ensemble mean
$\overline{\mathbf y}_b=M^{-1}\sum_{m=1}^{M}
\widehat{\mathbf y}_b^{(m)}$ as the point forecast:
\begin{equation}
  \operatorname{RMSE}
  =
  \left[
    \frac{1}{BD}\sum_{b=1}^{B}
    \left\|\overline{\mathbf y}_b-\mathbf y_b\right\|_2^2
  \right]^{1/2},
  \qquad
  \operatorname{MAE}
  =
  \frac{1}{BD}\sum_{b=1}^{B}
  \left\|\overline{\mathbf y}_b-\mathbf y_b\right\|_1.
  \label{eq:rmse-mae}
\end{equation}

\subsubsection{Stochastic ground truth}

For stochastic systems, the simulator provides a reference ensemble
$\{\mathbf y_b^{(r)}\}_{r=1}^{R}$ from $P_b$, which we compare with the model
ensemble from $Q_{\theta,b}$.  Before evaluation, each state variable is
standardized for each condition using the mean and population standard
deviation of the reference ensemble, pooled over its members and forecast
times.  The same transformation is applied to the model ensemble.  The
notation below refers to these standardized, vectorized trajectories.

\paragraph{Conditional mean error (CME) and conditional standard-deviation
error (CSE).}
Let $\boldsymbol{\mu}_{P,b}$ and $\boldsymbol{\mu}_{Q,b}$ denote the conditional
mean vectors under $P_b$ and $Q_{\theta,b}$, and let
$\boldsymbol{\sigma}_{P,b}$ and $\boldsymbol{\sigma}_{Q,b}$ denote their
coordinatewise conditional standard-deviation vectors.  We report the
relative errors
\begin{equation}
  \operatorname{CME}
  =
  \frac{1}{B}\sum_{b=1}^{B}
  \frac{\left\|\boldsymbol{\mu}_{P,b}-\boldsymbol{\mu}_{Q,b}\right\|_2}
       {\left\|\boldsymbol{\mu}_{P,b}\right\|_2},
  \qquad
  \operatorname{CSE}
  =
  \frac{1}{B}\sum_{b=1}^{B}
  \frac{\left\|\boldsymbol{\sigma}_{P,b}-\boldsymbol{\sigma}_{Q,b}\right\|_2}
       {\left\|\boldsymbol{\sigma}_{P,b}\right\|_2}.
  \label{eq:cme-cse}
\end{equation}
The moments are estimated from the respective ensembles using population
standard deviations.  In implementation, each denominator is lower-bounded
by a small positive constant for numerical stability.

\paragraph{First-order Wasserstein distance ($W_1$).}
Following its standard use in generative modeling
\citep{arjovsky2017wasserstein}, the population distance is
\begin{equation}
  W_1(P_b,Q_{\theta,b})
  =
  \inf_{\gamma\in\Pi(P_b,Q_{\theta,b})}
  \mathbb{E}_{(\mathbf Y,\widehat{\mathbf Y})\sim\gamma}
  \left[\left\|\mathbf Y-\widehat{\mathbf Y}\right\|_2\right],
  \label{eq:wasserstein-one}
\end{equation}
where $\Pi(P_b,Q_{\theta,b})$ is the set of couplings with marginals $P_b$
and $Q_{\theta,b}$.  For the two empirical ensembles, define the cost matrix
$C^{(b)}_{rm}=\|\mathbf y_b^{(r)}-\widehat{\mathbf y}_b^{(m)}\|_2$.
We compute the entropically regularized transport plan
\citep{cuturi2013sinkhorn}
\begin{equation}
  \gamma_{b,\varepsilon}^{\star}
  \in
  \argmin_{\gamma\in\mathcal{U}_{R,M}}
  \left\langle\gamma,C^{(b)}\right\rangle
  -\varepsilon H(\gamma),
  \qquad
  \widehat W_{1,\varepsilon}
  =
  \frac{1}{B}\sum_{b=1}^{B}
  \left\langle\gamma_{b,\varepsilon}^{\star},C^{(b)}\right\rangle,
  \label{eq:sinkhorn-wasserstein}
\end{equation}
where $\mathcal{U}_{R,M}$ contains nonnegative transport plans with uniform
marginals and $H$ is entropy.  We use $\varepsilon=0.1$ and 100 Sinkhorn
iterations.  Thus, the reported value is a Sinkhorn approximation to $W_1$,
not the debiased Sinkhorn divergence.

\paragraph{Energy distance (ED).}
For independent $\mathbf Y,\mathbf Y'\sim P_b$ and
$\widehat{\mathbf Y},\widehat{\mathbf Y}'\sim Q_{\theta,b}$, energy distance is
defined by \citep{szekely2013energy}
\begin{equation}
  \operatorname{ED}(P_b,Q_{\theta,b})
  =
  2\,\mathbb{E}\left\|\mathbf Y-\widehat{\mathbf Y}\right\|_2
  -\mathbb{E}\left\|\mathbf Y-\mathbf Y'\right\|_2
  -\mathbb{E}\left\|\widehat{\mathbf Y}-\widehat{\mathbf Y}'\right\|_2.
  \label{eq:energy-distance}
\end{equation}
We estimate this quantity with the unbiased finite-ensemble statistic
\begin{align}
  \widehat{\operatorname{ED}}_b
  ={}&
  \frac{2}{RM}\sum_{r=1}^{R}\sum_{m=1}^{M}
  \left\|\mathbf y_b^{(r)}-\widehat{\mathbf y}_b^{(m)}\right\|_2
  \nonumber\\
  &-
  \frac{1}{R(R-1)}
  \sum_{\substack{r,r'=1\\r\neq r'}}^{R}
  \left\|\mathbf y_b^{(r)}-\mathbf y_b^{(r')}\right\|_2
  \nonumber\\
  &-
  \frac{1}{M(M-1)}
  \sum_{\substack{m,m'=1\\m\neq m'}}^{M}
  \left\|\widehat{\mathbf y}_b^{(m)}-
  \widehat{\mathbf y}_b^{(m')}\right\|_2,
  \label{eq:energy-distance-ensemble}
\end{align}
and report
$\widehat{\operatorname{ED}}=B^{-1}\sum_{b=1}^{B}
\widehat{\operatorname{ED}}_b$.
Unlike the coordinatewise CRPS, both $W_1$ and ED use Euclidean distances
between complete trajectories and therefore evaluate the joint temporal
distribution.
